%------------------------------------- Ejercicio 1 -----------------------------------
\begin{ejer}
Sea $V$ un espacio vectorial de dimensión finita. Probar que si $\norm{\cdot}_1$ y $\norm{\cdot}_2$ son dos normas en $V$, entonces existe una constante $C \en \mathbb{R}$ tal que para todo $v \en V$ se cumple $\norm{v}_1 \leq C \norm{v}_2$.
\end{ejer}
\begin{sol}[Demostración]
Sea $\mathcal{B} = \set{b_1,\dots,b_n}$ una base para $V$, esta base induce un isomorfismo $T \colon V \to \mathbb{F}^n$. Luego, consideremos la norma $\norm{\cdot}'$ en $V$ inducida mediante $T$ por $(\mathbb{F}^n, \norm{\cdot})$, donde $\norm{\cdot}$ es la norma del taxista, es decir, se define como $\norm{v}' \coloneq \norm{T(v)}$; por lo tanto, si $v = \sum^n_{i=1}v_ib_i$, entonces 
\[
\norm{v}' = \norm{T(v)} = \norm{v_1,\dots,v_n} = \sum^n_{i=1} \abs{v_i}
\]
Con esta norma inducida, $T \colon (V,\norm{\cdot}') \to (\mathbb{F}^n,\norm{\cdot})$ es un isomorfismo isométrico. Luego, definimos
\[
M_1 \coloneq \max_{\mathclap{i=1,\dots,n}} \norm{b_i}_1%
\quad\text{y}\quad
M_2 \coloneq \max_{\mathclap{i=1,\dots,n}} \norm{b_i}_2%
\]
Entonces se cumple que, para cada $v \en V$ con $v = \sum^n_{i=1} v_ib_i$,
\[
\norm{v}_k%
\leq \norm\Big{\sum^{n}_{i=1} v_ib_i}{\vphantom{\sum}}_k%
\leq \sum^n_{i=1} \abs{v_i}\cdot\norm{b_i}_k%
\leq M_k \sum^n_{i=1} \abs{v_i}%
= M_k \norm{v}',
\quad\forall k\en\set{1,2}
\]
Esto prueba que $\norm{\cdot}_1,\norm{\cdot}_2 \colon (V,\norm{\cdot}') \to (\mathbb{R},\abs{\cdot})$ son Lipschitz, ya que se tiene para cada $v_1,v_2 \en V$
\[
\abs\big{\norm{v_1}_k - \norm{v_2}_k}%
\leq \norm{v_1 - v_2}_k%
\leq M_k \norm{v_1-v_2}',%
\quad\forall k\en\set{1,2}
\]
en particular, son continuas. Por otro lado, sabemos que $S \coloneq \set*{x \en \mathbb{F}^n \given \norm{x} = 1}$ es compacto en $(\mathbb{F}^n,\norm{\cdot})$ así que ---gracias a que $T^{-1}$ es continua--- $S' \coloneq T^{-1}(S)$ también es compacto en $(V, \norm{\cdot}')$ (nótese que $S' = \set{v \en V \given \norm{v}'4 = 1}$); de esto se sigue, por la continuidad de $\norm{\cdot}_2$, que existe $m \coloneq \min\set*{\norm{v}_2 \given v \en S'} > 0$. Ahora bien, esto implica que si $v \en V-\set0$ entonces 
\[
\hat{v} \coloneq \frac{v}{\norm{v}'} \en S'%
\quad\implies\quad%
m \leq \norm{\hat{v}}_2%
\quad\implies\quad%
m\norm{v}' \leq \norm{\hat{v}}_2\norm{v}' = \norm{v}_2
\]
Esto también se satisface para $v=0$, así que $m\norm{v}' \leq \norm{v}_2$ para todo $v \en V$. Finalmente, se tiene
\[
\norm{v}_1 \leq M_1\norm{v}' = \frac{M_1}{m}(m\norm{v}') \leq \frac{M_1}{m} \norm{v}_2,
\quad\forall v \en V
\]
Esto completa la prueba tomando $C = \frac{M_1}{m}$.
\end{sol}

%------------------------------------- Ejercicio 2 -----------------------------------
\begin{ejer}
Si $\mathcal{M}_n = \mathcal{M}_n(\mathbb{R})$ es el espacio vectorial (de dimensión $n^2$) de las matrices $n \times n$ a coeficientes reales, probar que $\norm{\cdot} \colon \mathcal{M}_n \to [0, +\infty)$ definida por
\[
\norm{A} = \sup \set*{ \norm{Av} \given v \en \mathbb{R}^n, \norm{v} = 1 }
\]
es una norma tal que para todo vector $v \en \mathbb{R}^n$ vale $\norm{Av} \leq \norm{A} \norm{v}$.
\end{ejer}
\begin{sol}[Demostración]
\emph{Primero veamos que está bien definida:}\\
Todo operador lineal $T \colon \mathbb{R}^n \to \mathbb{R}^n$ es continuo, en particular lo es $v \mapsto Av$; pero el conjunto $S \coloneq \set{x \en \mathbb{R}^n \given \norm{x} = 1}$ es compacto, por lo tanto, existe
\[
\max \set*{\norm{Av} \given v \en S} = \sup \set*{ \norm{Av} \given v \en \mathbb{R}^n, \norm{v} = 1 }
\]
\emph{Veamos que es una norma}:
\begin{itemize}
\item Como $\norm{Av} \geq 0$ para cualquier $v \en S$, entonces $\norm{A} \geq \norm{Av} \geq 0$, además, $\norm{0v} = 0$ para cada $v \en S$ así que $\norm{0} = 0$. Más aún, si $A \neq 0$ entonces existe $v \en \mathbb{R}^n$ con $Av \neq 0$ por lo que $v\neq0$ y $\norm{Av},\norm{v} > 0$. Luego, $\hat{v} \coloneq \frac{v}{\norm{v}} \en S$, de donde
\[
\norm{A} \geq \norm{A\hat{v}} = \frac{\norm{Av}}{\norm{v}} > 0
\]
Esto prueba que $\norm{\cdot}$ es \emph{definida positiva}.
\item Sean $A \en \mathcal{M}_n$ y $\lambda \en \mathbb{R}$. Como $\abs{\lambda} \geq 0$, es evidente que
\[
\norm{\lambda A} = \sup\set*{\norm{\lambda Av} \given v \en S}%
= \sup\set*{\abs{\lambda} \cdot \norm{Av} \given v \en S}%
= \abs{\lambda}\sup\set*{\norm{Av} \given v \en S}%
= \abs{\lambda}\cdot\norm{A}
\]
Esto prueba que $\norm{\cdot}$ es \emph{homogénea}.

\item Sean $A,B \en \mathcal{M}_n$. Entonces para cada $v \en S$ se cumple
\[
\norm{(A+B)v} = \norm{Av + Bv} \leq \norm{Av} + \norm{Bv} \leq \norm{A} + \norm{B}
\]
Esto muestra que $\norm{A} + \norm{B}$ es cota superior de $\set{\norm{(A+B)v} \given v \en S}$, por lo cual la definición de supremo implica que $\norm{A+B} \leq \norm{A} + \norm{B}$, i.e., $\norm{\cdot}$ satisface la \emph{desigualdad del triángulo}.
\end{itemize}\vspace{1em}
\emph{Veamos que $\norm{Av} \leq \norm{A}\cdot\norm{v}$}:\\
Si $v = 0$, ya terminamos, pues $Av = 0$ y $\norm{Av} = 0 \leq \norm{A}\cdot\norm{v}$. Supongamos que $v \en V - \set0$, entonces se tiene que $\norm{v} > 0$ y $\hat{v} \coloneq \frac{v}{\norm{v}} \en S$ por lo que
\[
\frac{\norm{Av}}{\norm{v}} = \norm{A\hat{v}} \leq \norm{A}%
\quad\implies\quad%
\norm{Av} \leq \norm{A} \cdot \norm{v}
\]
Esto termina la prueba.
\end{sol}

%------------------------------------- Ejercicio 3 -----------------------------------
\begin{ejer}
Probar que
\[
\norm{A}_0 = \sup\set*{ \abs*{a_{ij}} \given 1 \leq i, j \leq n }
\]
es otra norma en $\mathcal{M}_n$.
\end{ejer}
\begin{sol}[Demostración]
Notemos que como el supremo es tomado sobre un conjunto finito coincide con el máximo. Ahora probemos que $\norm{\cdot}_0$ es una norma sobre $\mathcal{M}_n$:
\begin{itemize}
\item Por definición, $\norm{A} \geq 0$ y como las entradas de la matriz $0$ son nulas, $\norm{0} = 0$. Más aún, si $A \neq 0$ entonces existen $i,j \en \set{1,\dots,n}$ tales que $a_{i.j} \neq 0$. Por lo tanto, $\norm{A}_0 \geq \abs{a_{i,j}} > 0$. Esto prueba que $\norm{\cdot}_0$ es \emph{definida positiva}.

\item Si $A \en \mathcal{M}_n$ y $\lambda \en \mathbb{R}$, la entrada $i,j\text{-ésima}$ de $\lambda A$ es $\lambda a_{i,j}$ por lo cual,
\[
\norm{\lambda A}_0%
= \sup_{\mathclap{1 \leq i, j \leq n}} \abs*{ \lambda a_{ij}}%
= \sup_{\mathclap{1 \leq i, j \leq n}} \abs{\lambda}\abs*{a_{ij}}%
= \abs{\lambda}\sup_{\mathclap{1 \leq i, j \leq n}} \abs*{a_{ij}}%
= \abs{\lambda}\cdot\norm{A}_0
\]
Esto prueba que $\norm{\cdot}_0$ es \emph{homogénea}.

\item Si $A,B \en \mathcal{M}_n$ entonces la $i,j\text{-ésima}$ entrada de $A+B$ es $a_{i,j} + b_{i,j}$  por lo que
\[
\abs{a_{i,j} + b_{i,j}} \leq \abs{a_{i,j}} + \abs{b_{i,j}} \leq \norm{A}_0 + \norm{B}_0,%
\quad\forall i,j \en \set{1,\dots,n}
\]
por lo que $\norm{A}_0 + \norm{B}_0$ es cota superior de $\set*{ \abs*{a_{ij} + b_{i,j}} \given 1 \leq i, j \leq n }$ así que la definición de supremo implica que $\norm{A+B}_0 \leq \norm{A}_0 + \norm{B}_0$. Esto prueba que $\norm{\cdot}_0$ satisface la \emph{desigualdad del triángulo}.
\end{itemize}
\end{sol}

%------------------------------------- Ejercicio 4 -----------------------------------
\begin{ejer}
Una función $A : \mathbb{R} \to \mathcal{M}_n$ es continua si para todo $t \en \mathbb{R}$ se tiene $\lim_{h\to 0} \norm{A(t+h) - A(t)} = 0$. Probar que $A(t)$ es continua si y solo si cada coeficiente $a_{ij}(t)$ es una función continua de $\mathbb{R}$ en $\mathbb{R}$.
\end{ejer}
\begin{sol}[Demostración]
\emph{Primero probémoslo para la norma $\norm{\cdot}_0$}:\\
Si $A(t)$ es continua, entonces para cada $i,j \en \set{1,\dots,n}$ y $t \en \mathbb{R}$
\[
0 \leq \abs{a_{i,j}(t+h) - a_{i,j}(t)} \leq \norm{A(t+h) - A(t)}_0
\]
por lo cual, gracias a que el lado derecho converge a $0$ cuando $h \to 0$, el teorema de compresión implica que $\lim_{h\to0} \abs{a_{i,j}(t+h) - a_{i,j}(t)} = 0$ para cada $t \en \mathbb{R}$, i.e., cada $a_{i,j}(t)$ es continua. Recíprocamente, supongamos que $a_{i,j}(t)$ es continua para cada $i,j \en \set{1,\dots,n}$ y sean $t \en \mathbb{R}$ y $\varepsilon > 0$. Por hipótesis, para cada $i,j \en \set{1,\dots,n}$ existe $\delta_{i,j} >0$ tal que 
\[
\forall h \en \mathbb{R}\;%
\Big(\;%
    \abs{h} < \delta_{i,j} \implies \abs{a_{i,j}(t+h) - a_{i,j}(t)} < \varepsilon
\;\Big)
\]
tomando $\delta \coloneq \min\set*{\delta_{i,j} \given 1\leq i,j\leq n}$ se cumple que
\[
\forall h \en \mathbb{R}\;\Big(\;%
\abs{h} < \delta \implies  \forall i,j \en \set{1,\dots,n},\; \abs{a_{i,j}(t+h) - a_{i,j}(t)} < \varepsilon
\;\Big)
\]
pero $\abs{a_{i,j}(t+h) - a_{i,j}(t)} < \varepsilon$ para cada $i,j \en \set{1,\dots,n}$ si y solo si $\norm{A(t+h) - A(t)} < \varepsilon$. Así que
\[
\forall h \en \mathbb{R}\;%
\Big(\;%
    \abs{h} < \delta \implies \norm{A(t+h) - A(t)} < \varepsilon
\;\Big)
\]
Como $\varepsilon > 0$ es arbitrario, esto prueba que $\lim_{h\to 0} \norm{A(t+h) - A(t)} = 0$ para cada $t \en \mathbb{R}$, es decir, $A(t)$ es continua.\par\vspace{1em}
\emph{Probémoslo para cualquier norma $\norm{\cdot}$}:\\
Por el \myref{ejer:4.1}, existen $C_1,C_2 > 0$ tales que $\norm{\cdot} \leq C_1 \norm{\cdot}_0$ y $\norm{\cdot}_0 \leq C_2 \norm{\cdot}$. Luego, 
\[
0 \leq \norm{A(t+h) - A(t)} \leq C_1\norm{A(t+h) - A(t)}_0%
\quad\text{y}\quad%
0 \leq \norm{A(t+h) - A(t)}_0 \leq C_2\norm{A(t+h) - A(t)}
\]
Si $A$ es continua respecto a $\norm{\cdot}_0$ entonces en la desigualdad izquierda $C_1\norm{A(t+h) - A(t)}_0 \xrightarrow{h\to0}0$ por lo que también $\norm{A(t+h) - A(t)} \xrightarrow{h\to0}0$ así que $A$ es continua respecto a $\norm{\cdot}$. Recíprocamente, si $A$ es continua respecto a $\norm{\cdot}$, en la desigualdad derecha $C_2\norm{A(t+h) - A(t)} \xrightarrow{h\to0}0$ por lo que también $\norm{A(t+h) - A(t)}_0 \xrightarrow{h\to0}0$ así que $A$ es continua respecto a $\norm{\cdot}_0$. Esto prueba que la continuidad respecto a ambas normas es equivalente, así que la prueba del ejercicio para $\norm{\cdot}_0$ sirve también como prueba del caso general con cualquier norma $\norm{\cdot}$.
\end{sol}

%------------------------------------- Ejercicio 5 -----------------------------------
\begin{ejer}
Sea $A:\mathbb{R} \to \mathcal{M}_n$ derivable, es decir, tal que cada coeficiente $a_{ij}(t)$ es una función derivable. Probar que la función $h(t) = \det(A(t))$ es derivable y que su derivada vale
\[
\dot{h}(t) = \sum_{k=1}^{n} \det(c^k_{ij}(t))
\]
donde $c^k_{ij}(t) = a_{ij}(t)$ si $i \neq k$ y $c^k_{ij}(t) = \dot{a}_{ij}(t)$ si $i = k$. En otras palabras, la matriz $(c^k_{ij})$ es la que se obtiene derivando la fila $k$ de la matriz $A$.
\end{ejer}
\begin{sol}[Demostración]
El determinante es definido como la única forma $n\text{-lineal}$ alternante (al verla como función de las filas o las columnas de la matriz) cuyo valor es $1$ en la matriz identidad. En este caso, la veremos como función de las filas de la matriz, i.e., si $A = (a_{i,j})^n_{i,j=1}$ entonces
\[
\det(a_{1,*}, a_{2,*}, \dots, a_{n,*}) \equiv \det A
\]
donde $a_{i,*}$ es la $i\text{-ésimo}$ fila de $A$.\par
\begin{lema}\label{lema:Binom}
Si $V \en \mathrm{FinVect}_{\mathbb{R}}$ y $D \colon V^k \to \mathbb{R}$ es una forma $k\text{-lineal}$, entonces para cada $(a_i),(b_i) \en V^k$,
\[
D(a_1 +b_1,\dots,a_{k}+b_{k})%
= \sum_{\mathclap{\substack{c_i \en \set{a_i,b_i}\\1\leq i \leq k}}} D(c_1,\dots,c_{k})
\]
\end{lema}
Este hecho se puede ver como una versión no conmutativa del producto de binomios y su prueba es una aplicación directa de inducción matemática sobre $k$. El caso base es trivial y el paso inductivo es el siguiente: si suponemos que la fórmula dada se cumple para cualquier forma $k\text{-lineal}$ y tomamos una forma $(k+1)\text{-lineal}$ $D \colon V^{k+1} \to \mathbb{R}$, entonces, dados $a_{k+1},b_{k+1} \en V$ tenemos que $D(-,-,\dots,a_{k+1})$ y $D(-,-,\dots,b_{k+1})$ son $k\text{-lineales}$, así que para cada $(a_i)^k_{i=1},(b_i)^k_{i=1} \en V^k$
\begin{align*}
D(a_1 +b_1,\dots,a_k+b_k,a_{k+1})%
&= \sum_{\mathclap{\substack{c_i \en \set{a_i,b_i}\\1\leq i \leq k}}} D(c_1,\dots,c_k,a_{k+1})\\
\\
&\;\text{y}\\
\\
D(a_1 +b_1,\dots,a_k+b_k,b_{k+1})
&= \sum_{\mathclap{\substack{c_i \en \set{a_i,b_i}\\1\leq i \leq k}}} D(c_1,\dots,c_k,b_{k+1})
\end{align*}
sumando ambas expresiones obtenemos
\[
D(a_1 +b_1,\dots,a_{k+1}+b_{k+1})%
= \sum_{\mathclap{\substack{c_i \en \set{a_i,b_i}\\1\leq i \leq k+1}}} D(c_1,\dots,c_{k+1})
\]
Esto completa la inducción.
\begin{coro}\label{coro:detcomb}
Si $A =(a_{i,j}), B = (b_{i,j}) \en \mathcal{M}_n$, entonces
\[
\det(A + B) = \sum_{C \en M} \det(C),%
\quad\text{donde}\quad%
M = \set*{(c_{i,j}) \en \mathcal{M}_n \given c_{i,*}= a_{i,*} \vee c_{i,*} = b_{i,*}}
\]
En particular, si $\lambda \en \mathbb{R}$,
\[
\det(A + \lambda B) = \det(A) + \lambda\sum_{\mathclap{C \en M_1}}\det(C) +\dots+ \lambda^n\sum_{\mathclap{C \en M_n}}\det(C)
\]
donde $M_k$ consta de aquellas matrices en $M$ que tienen $k$ filas provenientes de $B$.
\end{coro}
La primera parte del corolario es una aplicación directa del \myref{lema:Binom} con $D = \det$. La segunda se obtiene usando el primer resultado y particionando $M$ según la cantidad de filas provenientes de $B$ en cada matriz $C \en M$. Debido a la $n\text{-linealidad}$ del determinante, cada fila proveniente de $B$ aporta un factor $\lambda$ al correspondiente determinante.\par\vspace{1em}
\emph{Ahora podemos probar el resultado}:\\
Fijemos un $t \en \mathbb{R}$ arbitrario. Como existe $\dot{A}(t)$ ---cuyas entradas son los $\dot{a}_{i,j}(t)$---, se sigue que hay una función $r\colon \mathbb{R} \to \mathcal{M}_n$ tal que
\[
\lim_{\Delta t \to 0} \frac{r(\Delta t)}{\Delta t} = 0%
\quad\text{y}\quad%
\forall\Delta t \en \mathbb{R},\;A(t+\Delta t) - A(t) = \Delta t \cdot \dot{A}(t) + r(\Delta t)
\]
por lo cual, si definimos la función $\ell(\Delta t) \coloneq \dot{A}(t) + \Delta t^{-1}\cdot r(\Delta t)$, se tiene que
\[
\lim_{\Delta t \to 0} \ell
(\Delta t) = \dot{A}(t)%
\quad\text{y}\quad%
A(t+\Delta t) = A(t) + \Delta t \cdot \ell(\Delta t)
\]
Así que, gracias al \myref{coro:detcomb},
\begin{align*}
h(t + \Delta t) - h(t)%
&= \det(A(t + \Delta t))- \det(A(t))%
= \det(A(t) + \Delta t \cdot \ell(\Delta t)))- \det(A(t))\\
&= \Delta t \sum_{\mathclap{C \en M_1}} \det(C(\Delta t)) + \dots + \Delta t^n \sum_{\mathclap{C \en M_n}} \det(C(\Delta t))
\end{align*}
donde $M_l$ es el conjunto de las funciones $C(\Delta t)$ de $\mathbb{R}$ en $\mathcal{M}_n$ tales que $c_{i,*} \en \set{a_{i,*}(t), \ell_{i,*}}$ (nótese que $c_{i,*}$ es la función de $\mathbb{R}$ en $\mathbb{R}^n$ tal que $c_{i,*}(\Delta t)$ es la $i\text{-ésima}$ fila de $C(\Delta t)$) y $c_{i,*} = \ell_{i,*}$ para exactamente $l$ valores de $i$. Como el determinante es continuo con respecto a sus entradas (ya que es un polinomio con respecto a ellas) tenemos que, si $C \en M_l$ para algún $l \en \set{1,\dots,n}$, entonces 
\[
\lim_{\Delta t \to 0} \det(C(\Delta t)) = \det(\lim_{\Delta t \to 0} C(\Delta t))
\]
pero $\lim_{\Delta t \to 0} C(\Delta t)$ existe gracias a que existe el límite de cada una de sus entradas, más aún, conserva iguales aquellas filas provenientes de $A$ (ya que son constantes resp. a $\Delta t$) y cambia aquellas provenientes de $\ell(\Delta t)$ por su límite, es decir, por $\dot{a}_{i,*}(t)$. En particular, si $C \en M_1$ entonces $\lim C(\Delta t)$ es igual a algún $(c^k_{i,j})$. Luego
\begin{align*}
\dot{h}(t)
&= \lim_{\Delta t \to 0} \frac{h(t + \Delta t) - h(t)}{\Delta t}%
= \lim_{\Delta t \to 0} \left(\;%
    \sum_{\mathclap{C \en M_1}}\!\det(C(\Delta t)) + \dots + \Delta t^{n-1}\!\sum_{\mathclap{C \en M_n}}\!\det(C(\Delta t))
\;\right)
% \\
% &= \sum_{C \en M_1} \det(\lim_{\Delta t \to 0} C(\Delta t)) +0 \cdot \sum_{C \en M_2} \det(\lim_{\Delta t \to 0} C(\Delta t)) + \cdots + 0^{n-1} \cdot \det(\lim_{\Delta t \to 0} C(\Delta t))
= \sum^n_{k = 1}\!\det(c^k_{i,j})
\end{align*}

\end{sol}

%------------------------------------- Ejercicio 6 -----------------------------------
\begin{ejer}[Fórmula de Liouville] 
Probar que si $\phi:\mathbb{R} \to \mathcal{M}_n$ es la matriz principal en $t=0$ de la ecuación lineal $\dot{x} = A(t)x$ (i.e. $\dot{\phi}(t) = A(t)\phi(t)$ para todo $t \en \mathbb{R}$ y $\phi(0) = I$) entonces para todo $t_0 \en \mathbb{R}$ vale que
\[
\det\phi(t) = \det\phi(t_0) \ e^{\int_{t_0}^{t} \mathrm{tr}A(s)\,ds}
\]
donde $\mathrm{tr}A$ es la traza (suma de los elementos de la diagonal) de la matriz $A$.
\end{ejer}
\begin{sol}[Demostración]
\begin{nota}
Como el determinante de una matriz coincide con el de su transpuesta, el \myref{ejer:4.5} se cumple también si consideramos las columnas en lugar de las filas (es decir, si consideramos que $c^k_{i,j}$ se obtiene derivando la $k\text{-ésima}$ columna). De esta forma, también es conveniente cambiar la convención para que el determinante esté en función de las columnas de la matriz.
\end{nota}
Un corolario de que $\det \colon (\mathbb{R}^n)^n \to \mathbb{R}$ es la única forma $n\text{-lineal}$ alternante que vale $1$ en la identidad, es que cualquier otra forma $n\text{-lineal}$ alternante $D \colon (\mathbb{R}^n)^n \to \mathbb{R}$ es de la forma $D(A) = D(I_n)\det(A)$. A partir de esto, se prueba el siguiente lema:
\begin{lema}\label{lema:traza}
Para cualesquiera $A \en \mathcal{M}_n$ y $(v_1,\dots,v_n) \en (\mathbb{R}^n)^n$, se vale que
\[
\sum^n_{k=1} \det(v_1,\dots,Av_k,\dots,v_n) = \tr(A)\det(v_1,\dots,v_n)
\]
\end{lema}
Para probar este resultado, basta verificar que $\sum \det(-,\dots,A-,\dots,-)$ es alternante (claramente es $n\text{-lineal}$ porque cada uno de sus sumandos lo es) y vale $\tr(A)$ en la identidad (i.e., en $(e_1,\dots,e_n)$). Esto es así: supongamos que $v_i = v_j$ con $i < j$. Entonces hay dos casos para $k$
\begin{enumerate}[label = (\roman*)]
\item Para $k \noten \set{i,j}$ se tiene $\det(v_1,\dots,Av_k,\dots,v_n) = 0$ ya que el determinante es alternante.

\item Para $k \en \set{i,j}$ los sumandos correspondientes son
\[
\det(\dots,Av_i,\dots v_j,\dots)%
\quad\text{y}\quad
\det(\dots,v_i,\dots Av_j,\dots)%
\]
pero como el determinante es alternante y $v_i = v_j$
\[
\det(\dots,v_i,\dots Av_j,\dots) = -\det(\dots,Av_j,\dots v_i,\dots) = -\det(\dots,Av_i,\dots v_j,\dots)
\]
por lo que estos dos sumandos se van a cancelar
\end{enumerate}
Sumando sobre $k \en \set{1,\dots,n}$ obtenemos que cuando $v_i = v_j$
\[
\sum^n_{k=1} \det(v_1,\dots,Av_k,\dots,v_n) = 0
\]
Esto prueba que $\sum\det(-,\dots,A-,\dots,-)$ es una forma $n\text{-lineal}$ alternante; además,
\[
\sum^n_{k=1} \det(e_1,\dots,Ae_k,\dots,e_n) = \sum^n_{k=1} \det(e_1,\dots,a_{*,k},\dots,e_n) = \sum^n_{k=1}a_{k,k} = \tr(A)
\]
Esto demuestra el lema.\clearpage
\emph{Ahora probemos el ejercicio}:\\
Por el \myref{ejer:4.5} aplicado a las columnas, la derivada de $h(t) \coloneq \det(\phi(t))$ es
\[
\dot{h}(t) = \sum^n_{k=1} \det(\phi_{*,1}(t),\dots,\dot{\phi}_{*,k}(t),\dots\phi_{*,n}(t))
\]
donde $\phi_{*,j}(t)$ denota la $j\text{-ésima}$ columna de $\phi(t)$ con $1 \leq j \leq n$. No obstante, sabemos que las columnas de la matriz fundamental son soluciones de la EDO, así que el \myref{lema:traza} implica
\[
\dot{h}(t) = \sum^n_{k=1} \det(\phi_{*,1}(t),\dots,A(t)\phi_{*,k}(t),\dots,\phi_{*,n}(t)) = \tr(A(t))\det(\phi_{*,1}(t),\dots,\phi_{*,n}(t)) = \tr(A(t))h(t)
\]
De esto se sigue que $h(t)$ es solución de la EDO lineal homogénea $\dot{x} = \tr(A(t))x$. Sin embargo, ya sabemos que la única solución a esta EDO para la condición inicial $(t_0,h(t_0))$ es
\[
h(t_0)e^{\int_{t_0}^{t} \mathrm{tr}A(s)\,ds}
\]
Por lo tanto, $h(t) = h(t_0)e^{\int_{t_0}^{t} \mathrm{tr}A(s)\,ds}$, es decir,
\[
\det\phi(t) = \det\phi(t_0) \ e^{\int_{t_0}^{t} \mathrm{tr}A(s)\,ds}
\]
\end{sol}

%------------------------------------- Ejercicio 7 -----------------------------------
\begin{ejer}
Halle los puntos de equilibrio del sistema
\[
\begin{cases}
\;\dot{x} = 2y(z-1) \\
\;\dot{y} = -x(z-1) \\
\;\dot{z} = -z^3
\end{cases}
\]
Muestre que la linealización alrededor del origen no permite decidir su estabilidad. Para determinar la estabilidad, busque una función de Liapunov de la forma \[V(x,y,z) = ax^2 + by^2 + cz^2\]
\end{ejer}
\begin{sol}
Si $f(x,y,z) = (0,0,0)$ entonces $-z^3 = 0$ por lo cual $z = 0$ así que sustituyendo: $-2y = 2y(z-1) = 0$ y $x = -x(z-1) = 0$, es decir, $(x,y,z) = (0,0,0)$. Se sigue que el único punto de equilibrio es el origen.

La matriz jacobiana del campo $f(x,y,z) = (2y(z-1), -x(z-1), -z^3)$ es
\[
J_f(x,y,z) =%
\mqty(%
0 & 2(z-1) & 2y\\
-(z-1) & 0 & -x\\
0 & 0 & -3z^2
)
\]
Así que evaluada en cero da $J_f(0,0,0) = \smqty(0&-2&0\\1&0&0\\0&0&0)$. Esta matriz tiene kernel no trivial (pues $e_3$ está en su kernel) lo cual implica que tiene a $0$ como valor propio así que no se puede determinar su estabilidad a partir de su linealización en el origen (ya que no todos sus valores propios tienen parte real negativa).

En su defecto, determinemos la estabilidad buscando una función de Liapunov de la forma 
\[
V(x,y,z) = ax^2 + by^2 + cz^2
\]
Calculamos la derivada sobre una solución hipotética 
\[
\dot{V}(x,y,z) = \ang{\nabla V \given f} = \ang{(2ax,2by,2cz) \given (2y(z-1),-x(z-1),-z^3)} = (4a-2b)xy(z-1) - 2cz^4
\]
como $xy(z-1)$ no mantiene el mismo signo en ningún entorno del origen, basta con que $4a-2b = 0$. Entonces podemos proponer
\[
V(x,y,z) = \frac14x^2 + \frac12y^2 + \frac12z^2
\]
Esta función es definida positiva y su derivada $\dot{V}(x,y,z) = -z^4$ es semidefinida negativa. La existencia de esta función de Liapunov implica que el origen es punto de equilibrio estable.
\end{sol}

%------------------------------------- Ejercicio 8 -----------------------------------
\begin{ejer}
Consideramos la ecuación diferencial autónoma $\dot{x} = F(x)$ donde $F$ es un campo de vectores acotado y localmente Lipschitz en $\mathbb{R}^n$, tal que $p \en \mathbb{R}^n$ es un punto de equilibrio asintóticamente estable. Denotamos $\varphi(t,x)$ al flujo en $\mathbb{R}^n$ que define esta ecuación, definimos la \textit{cuenca de atracción} del punto $p$ como el conjunto
\[
A_p = \set[\big]{ x \en \mathbb{R}^n \given \lim_{t \to \infty} \varphi(t,x) = p }.
\]
\begin{enumerate}[label=(\alph*)]
\item Probar que $A_p$ es abierto.

\item Probar que tanto $A_p$ como su borde $\partial A_p$ son conjuntos invariantes por el flujo.

\item Probar que todo punto $z \en \partial A_p$ es inestable y concluir que si todos los puntos son estables para la ecuación, entonces todas las órbitas tienden a $p$ en el futuro.
\end{enumerate}
\end{ejer}
\begin{sol}[Demostración]
Por practicidad, primero probaré el inciso (b):
\begin{enumerate}
\item[(b)] \emph{Veamos que $A_p$ es invariante por la acción del flujo}:\\
Por resultados elementales sobre límites y ya que $\varphi(t+t_0,x) = \varphi(t,\varphi(t_0,x))$, se cumple que para cualquier $t_0 \en \mathbb{R}$
\[
\lim_{t\to\infty}\varphi(t,x) = p%
\quad\iff\quad%
\lim_{t\to\infty}\varphi(t + t_0,x) = p%
\quad\iff\quad%
\lim_{t\to\infty}\varphi(t, \varphi(t_0,x)) = p
\]
Es decir, $x \en A_p$ si y solo si $\varphi(t_0,x) \en A_p$. Como $t_0 \en \mathbb{R}$ es arbitrario, esto demuestra que $A_p$ es invariante bajo el flujo.\footnote{Más aún, probamos que también el complemento es invariante bajo el flujo: como es un "si y solo si", en particular $x \noten A_p$ implica que para cada $t_0 \en \mathbb{R}$ se tiene $\varphi(t_0,x) \noten A_p$.}\par
\emph{Veamos que $\partial A_p$ es invariante por la acción del flujo}:\\
Sean ahora $x \en \partial A_p$ y $t_0 \en \mathbb{R}$ arbitrarios. Entonces, dado $\varepsilon > 0$ arbitrario, tenemos (nuevamente, por continuidad respecto a condiciones iniciales) que existe $\delta > 0$ tal que
\[
\norm{x - y} < \delta%
\quad\implies\quad%
\forall t \en I,\; \norm{\varphi(t,x) - \varphi(t,y)} < \varepsilon
\]
donde $I = [\min\set{0,t_0},\max\set{0,t_0}]$. Sin embargo, como $x \en \partial A_p$ se sigue que existen $y,z \en B(x,\delta)$ tales que $y \en A_p$ y $z \noten A_p$. Esto implica que $\varphi(t_0,y), \varphi(t_0,z) \en B(\varphi(t_0,x), \varepsilon)$ y (por lo visto en el inciso anterior) $\varphi(t_0,y) \en A_p$ y $\varphi(t_0,z) \noten A_p$. Como $\varepsilon>0$ es arbitrario, concluimos que $\varphi(t_0,x) \en \partial A_p$.

\item[(a)] Sea $x \en A_p$. Como $p$ es asintóticamente estable, existe $\delta_1 > 0$ tal que para cada $y \en \mathbb{R}^n$
\[
\norm{y-p} < \delta_1%
\quad\implies\quad%
\lim_{t \to \infty} \varphi(t,y) = p%
\quad\text{(equivalentemente, $y \en A_p$)}
\]
es decir, $B(p,\delta_1) \subseteq A_p$. Además, como $x \en A_p$, existe $t_1 \geq 0$ tal que
\[
\forall t\geq t_1, \quad\norm{\varphi(t,x) - p} < \frac{\delta_1}{2}
\]
Por otro lado, la continuidad respecto a las condiciones iniciales implica que existe $\delta_2 > 0$ tal que
\[
\norm{y - x} < \delta_2%
\quad\implies\quad%
\forall t\en[0,t_1],\;\norm{\varphi(t,y) - \varphi(t,x)} < \frac{\delta_1}{2}
\]
De esto se sigue que para cada $y \en B(x, \delta_2)$ se satisface
\[
\norm{\varphi(t_1,y) - p} \leq \norm{\varphi(t_1,y) - \varphi(t_1,x)} + \norm{\varphi(t_1,x) - p} < \frac{\delta_1}{2} + \frac{\delta_1}{2} = \delta_1
\]
lo cual implica $\varphi(t_1,y) \en A_p$
así que, por lo visto en (b), $y \en A_p$. Hemos probado que $B(x,\delta_2) \subseteq A_p$ y, como $x \en A_p$ es arbitrario, se sigue que $A_p$ es abierto.

\item[(c)] Supongamos por contradicción que existe $z \en \partial A_p$ estable. Por un lado, la estabilidad asintótica de $p$ garantiza que existe $\delta_1 >0$ tal que para cada $y \en \mathbb{R}^n$
\[
\norm{y-p} < \delta_1%
\quad\implies\quad%
\lim_{t\to\infty} \varphi(t,y) = p%
\quad\text{(equivalentemente, $y \en A_p$)}
\]
Por otro lado, la estabilidad de $z$ implica que existe $\delta_2 > 0$ tal que para cada $y \en \mathbb{R}^n$
\[
\norm{z-y} < \delta_2%
\quad\implies\quad%
\forall t\geq0,\;\norm{\varphi(t,z) - \varphi(t,y)} < \frac{\delta_1}{2}
\]
Sin embargo, como $z \en \partial A_p$ se sigue que existe $y_0 \en B(z,\delta_2) \cap A_p$; en particular, como $y_0 \en A_p$ se tiene que $\lim_{t\to\infty}\varphi(t,y_0) = p$ así que existe $t_0\geq0$ tal que
\[
\forall t \geq t_0,\;\norm{\varphi(t,y_0) - p} < \frac{\delta_1}{2}
\]
Finalmente, como también tenemos $\norm{z-y_0} < \delta_2$ entonces se cumple
\[
\norm{\varphi(t_0,z) - p} \leq \norm{\varphi(t_0,z) - \varphi(t_0,y_0)} + \norm{\varphi(t_0,y_0) - p} < \frac{\delta_1}{2} + \frac{\delta_1}{2} = \delta_1
\]
esto implica que $\varphi(t_0,z) \en A_p$ así que (por inciso (b)) $z \en A_p$. No obstante, esto implica que $z \en \partial A_p \cap A_p$, lo cual es una contradicción, pues ya vimos que $A_p$ es abierto por lo que $\partial A_p \cap A_p = \varnothing$. El error estuvo en suponer que $z \en \partial A_p$ es estable, por lo que todos los puntos en la frontera son inestables.\par
En particular, si todos los puntos son estables se sigue que $\partial A_p = \varnothing$, de donde se sigue que $A_p$ es un abierto sin frontera de $\mathbb{R}^n$, como $\mathbb{R}^n$ es conexo, esto implica que $A_p = \mathbb{R}^n$. Es decir, todas las órbitas convergen a $p$ en el futuro.
\end{enumerate}
\end{sol}